\section{The Propers: Detailed Commentary}

\subsection{Confession inside the furnace (Introit)}
Azarias prays from within the furnace. His companions have refused the king's
idolatrous command; their fidelity has brought them into danger. Yet his
prayer speaks in the first-person plural of sin, disobedience, judgment,
and mercy. The Introit selects and rearranges this language, joining it to
the opening verse of Psalm~118. The result moves from an acknowledgment
of disobedience toward the happiness of walking in God's law.

Jerome confronts the apparent contradiction: the three young men are not
being punished for the offences they confess. At Daniel~3:29 he explains
that they speak in the person of the people. Their voice of solidarity is
larger than a private inventory of their own acts. This makes confession
possible without making the faithful witnesses guilty of the king's
idolatry.\footnote{Jerome, \work{Commentary on Daniel}, I, at 3:23--29,
especially the note on 3:29, PL~25, cols.~639--640.}
The assembly likewise acknowledges its need for mercy together. Nothing
in this confession identifies a particular illness as the punishment of a
particular sufferer.

The wider chapter makes the prayer's trust more demanding. Before the
furnace, the young men confess God's power to deliver and refuse idolatry
even if deliverance does not come (Dan.~3:17--18). Their faith is not a
bargain with a guaranteed visible result. Nor is Psalm~118 a celebration
of unaided moral achievement. Its praise of the undefiled way belongs to
a poem that repeatedly asks to be taught, helped, forgiven, and restored.
Augustine reads its opening beatitude with this dependence on divine
help.\footnote{Augustine, \work{Expositions of the Psalms}, Ps.~118,
\S\S1--8; NPNF I.8, heading Ps.~119.}
The entry into worship therefore unites responsibility and need. Sin is
confessed rather than excused, and mercy is sought rather than presumed.

\subsection{Pardon that makes service possible (Collect)}
The Collect asks for two gifts together, pardon and peace, and states their
purpose: cleansed servants may serve God with a secure mind. Its peace is
active. The mind is freed for service, not merely relieved of an unpleasant
feeling. Schuster explains the relation by describing a heart divided by
its passions and remorse: forgiveness restores the possibility of peaceful
attention and good action.\footnote{Schuster, \work{The Sacramentary}, III,
``Twentieth Sunday after Pentecost,'' p.~175.}
This prayer gives the Introit's confession a destination. The same people
who acknowledge disobedience ask to become capable of faithful service.

When the Maternity Collect follows, the petition names the Word who truly
took flesh in Mary's womb and asks her intercession. Divine power and
human birth meet in one person. Augustine ends his sixteenth tractate on
John by contemplating the same wonder: the Maker of Zion was made man
within it, receiving flesh from the Virgin. His argument illuminates the
prayer's Christology. The request for Mary's help rests in the incarnate
Son's saving work; it does not introduce another source of divine power.
\footnote{Augustine, \work{Tractates on John}, XVI.7; Maternity Collect,
\work{Missale Romanum} (1962), p.~685, n.~3847.}

\subsection{Wisdom that becomes mutual service (Epistle)}
The Epistle's first verb concerns walking. Ephesians has already described
the transformed life in truthfulness, forgiveness, and love; now the
assembly hears how carefully that life must be lived. Redeeming time
belongs to discernment of God's will. It is more than working faster or
making each hour productive. Chrysostom places the command amid hostility
and the danger of retaliation: Christians must avoid creating needless
occasions of offence while refusing to surrender the Gospel. They may have
to relinquish advantages in order to keep their fidelity.
\footnote{Chrysostom, \work{Homilies on Ephesians}, XIX, opening exposition
of 5:15--17; NPNF I.13.}

The contrast between wine and the Spirit then concerns the kind of life
that fills a person and shapes a community. Drunkenness dissolves
attentiveness; spiritual song demands it. Chrysostom explains singing in
the heart as singing with understanding, so that the mouth and the inward
attention do not travel in different directions. The words are addressed
to one another as well as to God. The end of the appointed reading is
therefore essential: mutual submission in the fear of Christ is one of
the forms this Spirit-filled life takes.

Chrysostom develops that reciprocity through a comparison. A household of
people compelled to serve one master differs from friends who willingly
serve one another. His appeal turns even the master toward service and
provision for those under him. The ancient social arrangement is present
in his example, but the exegetical pressure falls on the person who would
claim exemption from serving. Worship does not authorize such exemption.
\footnote{Chrysostom, \work{Homilies on Ephesians}, XIX, exposition of
5:18--21 and concluding moral exhortation.}

\subsection{The open hand and the prepared heart (Gradual and Alleluia)}
Psalm~144 praises God's kingdom, mercy, faithfulness, and care for living
creatures. The Gradual's universal eyes and open hand belong first to
this praise of the Creator. Food arrives in its fitting season. Augustine
compares the giver to a physician who knows both what a patient should
receive and when. His point is to sustain a petitioner whose good request
has not received an immediate answer. In the following verses he directs
desire beyond benefits to God himself.\footnote{Augustine, \work{Expositions
of the Psalms}, Ps.~144, \S\S14--18; NPNF heading Ps.~145.}

Bellarmine brings out a different demand in the same universal language.
The creatures' dependence on God's generosity exposes human hoarding and
misuse of resources. Those who rule and possess must imitate the open hand.
Schuster, considering this chant within the Mass, develops its Eucharistic
meaning: nourishment given under the veil of faith leads toward the sight
of God. The Psalm's creaturely food, the moral duty to share, and this
sacramental fulfilment give the chant breadth. None requires treating
ordinary bodily hunger as unimportant.\footnote{Bellarmine, \work{Commentary
on the Psalms}, Ps.~144:14--17; Schuster, III, pp.~175--176.}

The Alleluia answers with a prepared heart. Its biblical source, Psalm~107,
also knows threatened circumstances and the insufficiency of merely human
help. Readiness is thus compatible with dependence. The chant's wording
differs from the biblical verse: it addresses \emph{tibi, gloria mea},
where the historical biblical English has ``with my glory.'' Schuster
hears God addressed as the soul's glory. That gives praise a further
freedom: human approval is not its final measure. Augustine's brief
exposition of Psalm~107 identifies its reuse of earlier psalms; it does
not itself supply Schuster's explanation of this liturgical address.
\footnote{Augustine, Ps.~107, \S\S1--2 (NPNF Ps.~108); Schuster, III, p.~176.}

\subsection{Two moments of belief (Gospel)}
John gives the encounter a precise narrative order. The ruler hears that
Jesus has returned to Galilee; he comes because his son is dying; he asks
Jesus to come down. Jesus' rebuke speaks in the plural of dependence on
signs. The father repeats his urgent request. Christ then sends him away
with assurance that the son lives, and the man believes the word and goes.
On the way the servants report the recovery. His inquiry about the hour
reveals its correspondence with Jesus' word, and the narrative ends with
his belief and that of his whole household.

The belief at verse~50 and the belief at verse~53 are not redundant.
Trust has already changed the father's action before the report, while
the report makes the healing's relation to Jesus manifest and the ending
includes the household. The narrative itself supplies these distinct
moments without giving a universal scheme of stages through which every
believer must pass. Cana, where Jesus speaks, is named in the preceding
part of verse~46; Capharnaum is where the child is ill. The distance matters
because the healing requires no bodily journey by Jesus, not because the
reader can calculate an exact mileage from the passage.

Augustine interprets the ruler's initial approach severely, as coldness
or absence of faith exposed by Christ. Gregory argues instead that asking
Christ to heal already shows faith, but an imperfect faith that mistakenly
requires bodily presence. Their estimates of the starting point differ.
Both recognize that the petitioner's understanding of Christ's power must
be transformed.\footnote{Augustine, \work{Tractates on John}, XVI.3--5;
Gregory, \work{Homilies on the Gospels}, II.28.1.}
The healing is real, and the later report is evidence. Christ's correction
of dependence on signs does not make the sign useless: the unappointed
verse~54 expressly calls the event a sign.

\subsection{Tears offered with the gifts (Offertory)}
Psalm~136 remembers Jerusalem from Babylon. The liturgical adaptation
turns remembrance into direct address to Zion. In the full psalm the
captors demand a song and the exiles resist turning sacred memory into
entertainment. Only the opening lament is appointed here, but its context
gives the tears their object: a home and a worship that have been violently
lost. The Offertory does not ask the assembly to invent an agreeable
feeling before it approaches the altar.

Augustine develops the contrast between the rivers that flow away and
Zion's abiding home. To sit beside the currents rather than plunge into
them is to acknowledge captivity without surrendering one's desire to
transitory goods. His conclusion calls the pilgrims to sing to one
another. Lament and faithful song belong to the same journey.
Bellarmine likewise distinguishes outward life amid worldly obligations
from inward captivity to them.\footnote{Augustine, Ps.~136, \S\S2--9,13
(NPNF Ps.~137); Bellarmine, Ps.~136:1--4.}
The violent final verse of the full psalm is outside the chant. Augustine
reads its little ones as evil desires to be overcome in their beginnings
by Christ; his spiritual exposition is directed against vice, not people.

\subsection{Medicine and the worship of the Trinity (Secret and Preface)}
The Secret asks that these mysteries provide heavenly medicine and expel
the vices of the heart. The plural vices gives moral healing specificity:
the heart's desires can be disordered and need cleansing. Schuster explains
the Eucharistic gift as remedy against sin, not simply as a reward displayed
before those already healthy.\footnote{Schuster, III, p.~177.}
Where the Maternity Secret follows, its petition adds present and perpetual
prosperity and peace through divine mercy and Mary's intercession. The
temporal and enduring gifts remain together; present need does not cease
to matter when eternal peace is desired.

The Trinity Preface gives thanksgiving its explicit confession. The Father,
his only-begotten Son, and the Holy Spirit are one God and one Lord; the
unity is not a solitary Person. Distinction of Persons, unity of essence,
and equality of majesty belong together. Ephesians has already directed
Spirit-filled thanks to the Father in the name of Jesus Christ. The
Preface's confession gives this thanksgiving its full liturgical address.
The angels' praise joins the assembly's song at a point beyond its changing
circumstances. Its object is God's majesty, not the assembly's success.

\subsection{A promise remembered and a life made obedient (Communion and Postcommunion)}
The Communion asks God to remember the word that gave his servant hope.
Augustine explains that God has not suffered a failure of memory. The
petition asks him to fulfil what he has promised; it manifests the servant's
desire. Humiliation includes the lowliness of penitence, trial, and
persecution, and the promise comforts the Church even when death threatens.
\footnote{Augustine, Ps.~118, \S\S50--56 (NPNF Ps.~119).}
The biblical verse continues by saying that the word has given life.
That clause supplies context but is omitted from the appointed chant.
Schuster brings the Communion's word into a Christological focus: the
incarnate Word is offered to the Father as the ground of hope.
\footnote{Schuster, III, p.~177.}

The Postcommunion asks for continual obedience so that the faithful may
be worthy of the sacred gifts. It does not announce that they have earned
those gifts. Obedience itself is requested from the Lord. Schuster relates
this habitual docility to fruitful sacramental reception; Augustine's
account of righteousness given through grace explains why humility belongs
to that docility. The prayer returns to the Introit's commandments with a
changed accent: disobedience has been confessed, and obedience is now
asked as a gift. In the commemorative branch the final Marian prayer adds
cleansing and fellowship in heavenly remedy. Lasance's English identifies
that remedy personally with Christ. This is the historical translators'
interpretive rendering of the Latin prayer, not an additional Latin title.
